\documentclass[oneside,spanish,10pt]{amsart}
\usepackage{setspace}
\usepackage[a4paper]{geometry}
\geometry{verbose,tmargin=2cm,bmargin=2cm,lmargin=2.5cm,rmargin=2.5cm,footskip=1cm,includefoot,headsep=1cm,includehead} % Tamaño
\usepackage{amsthm}
\usepackage[spanish]{babel} % Idioma español
\usepackage[shortlabels]{enumitem} % Mejora el entorno enumerate
\usepackage{graphicx} % Colocar imágenes
\usepackage[hidelinks, colorlinks=true, citecolor=black, linkcolor=., urlcolor=blue]{hyperref} % Vínculos y referencias interactivas
\usepackage{xurl}
\usepackage[skip=3pt]{caption} % Permite modificar el espacio entre el caption y la imagen
\usepackage{multicol} % Entornos de múltiples columnas
\usepackage{nicefrac}
\usepackage{tikz}
\usepackage[center]{titlesec}
\usepackage{titletoc}
\usepackage{pdfpages}
\usepackage{tcolorbox}

%-------------------------------------------------------------------------
% Configuraciones iniciales
\makeatletter
%Numeración
\numberwithin{equation}{section}
\numberwithin{figure}{section}
\newlength{\lyxlabelwidth}      % Longitud auxiliar
\makeatother
%-------------------------------------------------------------------------

%-------------------------------------------------------------------------
% FUENTES
\usepackage[medium]{inter}
\usepackage{mathpazo}
\usepackage{domitian}
\usepackage[T1]{fontenc}
\let\oldstylenums\oldstyle
%-------------------------------------------------------------------------

%-------------------------------------------------------------------------
% Bibliografía
\usepackage[%
backend=biber,
bibstyle=apa,
citestyle=numeric,
sorting=none,
block=none,
maxbibnames=99
]{biblatex}

\defbibheading{bibliography}[\refname]{}

\renewcommand*{\labelalphaothers}{}

\makeatletter
\RequireBibliographyStyle{numeric}
\makeatother

%-------------------------------------------------------------------------
% ARCHIVOS DE BIBLIOGRAFÍAS
\addbibresource{Trabajos/01/refs-01.bib}
\addbibresource{Trabajos/02/refs-02.bib}
\addbibresource{Trabajos/03/refs-03.bib}
%-------------------------------------------------------------------------

%-------------------------------------------------------------------------
% Otras configuraciones
\addto\captionsspanish{\renewcommand{\tablename}{Tabla}}
\addto\captionsspanish{\renewcommand{\figurename}{Figura}}

\theoremstyle{definition}
\newtheorem{actividad_previa}{\normalfont ACTIVIDAD PREVIA}
\newtheorem{actividad}{\normalfont ACTIVIDAD}
\newtheorem{actividad_entre}{\normalfont ACTIVIDAD ENTRE CLASES}
\newtheorem*{evaluacion}{\normalfont EVALUACIÓN}
\newtheorem*{observacion}{Observacion}
\newtheorem*{teorema}{Teorema}
%-------------------------------------------------------------------------

\usepackage{fancyhdr}
\pagestyle{fancy}
\renewcommand{\headrulewidth}{0pt}
\fancyhead[C]{\scshape Memorias de las VIII JEM - 2024\\ \vspace{1mm} \hrule \vspace{1mm}
	\scshape\scriptsize \articulotitulocorto}
\fancyhead[L]{}
\fancyhead[R]{\vspace{1em}\small\scshape\articulotipo}
\fancyfoot[C]{\small \thepage}

\fancypagestyle{portadapage}{
	\fancyhf{}
	\renewcommand{\headrulewidth}{0pt}
	\fancyhead[C]{\scshape Memorias de las VIII JEM - 2024}
	\fancyhead[R]{\small\scshape\articulotipo}
	\fancyfoot[C]{\small \thepage}
}

%-------------------------------------------------------------------------
% TABLA DE CONTENIDOS
\setcounter{tocdepth}{1}
\newcommand{\desctotoc}[1]{%
	\addtocontents{toc}{\hspace*{3em}\medskip\detokenize{\textit{#1}}\leavevmode\par\medskip}
}
%-------------------------------------------------------------------------

%-------------------------------------------------------------------------
% TÍTULOS DE SECCIONES
\setcounter{secnumdepth}{4}
\titleformat{\section}[block]{\centering\setstretch{0.75}\Huge\scshape}{}{1em}{}
\titleformat{\subsection}{\centering\scshape}{\arabic{subsection}.}{1em}{\LARGE}
\titleformat{\subsubsection}{\normalfont\bfseries}{\arabic{subsection}.\arabic{subsubsection}}{1em}{\large}
\titleformat{\paragraph}{\normalfont\bfseries}{\arabic{subsection}.\arabic{subsubsection}.\arabic{paragraph}}{1em}{\large}
\titlespacing*{\paragraph}{0pt}{1em}{1em}

\titlecontents{section}
[3em] % indentación desde margen izquierdo
{\vspace{2pt}} % formato antes de la entrada (puedes dejar vacío {})
{\contentslabel{1.5em}} % formato del número (ancho reservado)
{} % formato sin número (si no tiene)
{\hfill\contentspage} % formato del número de página alineado a la derecha


%-------------------------------------------------------------------------

%-------------------------------------------------------------------------
% QUOTING
\let\oldquote\quote \renewcommand\quote{\oldquote\small}
%-------------------------------------------------------------------------

%-------------------------------------------------------------------------
% REFERENCIAS
\renewcommand{\thesubsubsection}{\arabic{subsection}.\arabic{subsubsection}}
%-------------------------------------------------------------------------

%-------------------------------------------------------------------------
% AUTORES
\newcommand{\addautor}[3][]{%
	{\Large \scshape #2}\\%
	{\large \itshape #3}\\%
	\ifx\\#1\\%
	% Si el argumento opcional (correo) está vacío, no hace nada
	\else
	{\ttfamily \small #1}\\%
	\fi
}
%-------------------------------------------------------------------------

%-------------------------------------------------------------------------
% DOCUMENTO

\title{}
\author{}
\date{}
\newcommand{\articulotitulo}{}
\newcommand{\articulotitulocorto}{}
\newcommand{\articulotipo}{}

\begin{document}
	
	% Toda la primer parte
	\includepdf{Trabajos/Portada.pdf}
	\clearpage
	\begin{center}
	\LARGE \bfseries \sffamily
	Memorias de las VI Jornadas de Enseñanza de la Matemática (JEM)
\end{center}

\vspace*{2.5mm}

\begin{center}
	\begin{tcolorbox}[colback=white, arc=0mm, width=10cm, boxrule=0.5pt]
		\footnotesize
		Universidad Nacional de Salta
		
		\hspace*{1em} Memorias de las VI Jornadas de Enseñanza de la Matemática (JEM) / Compilado por Sángari, Antonio Noé - Salta: Universidad Nacional de Salta, 2023.
		
		\vspace*{1em}
		
		\hspace*{1em} Archivo digital: descarga y online
		
		\hspace*{1em} ISBN 978-987-633-615-4		
	\end{tcolorbox}
\end{center}

\vspace*{2.5mm}

\begin{center}
	\sffamily \Large \bfseries Compiladores
	
	\large \normalfont
	Guzmán González, Ramiro 
	
	Sángari, Antonio Noé 
\end{center}

\begin{center}
	\sffamily \Large \bfseries Revisores de talleres y comunicaciones breves
	
	\vspace*{-1em}
	
	\begin{multicols}{2}
		\large \normalfont
		Alberto, Diego Luis
		
		Aliendro, Estela Sonia
		
		Arias, Jesús Eduardo
		
		Ávila, Mario Ubaldo
		
		Bifano, Fernando Jorge  
		
		Duarte, Betina
		
		Etchegaray, Silvia Catalina
		
		Ferrero, María Martha 
		
		Formeliano, Blanca Azucena 
		
		Funes, Héctor Nicolás
		
		García, José Ignacio
		
		Martínez, Rosa
		
		Marrón, Beatriz Susana
		
		Méndez, Nilda Graciela
		
		Mercau, Beatriz Susana
		
		Oliva, Elisa
		
		Quintana, Pablo
		
		Rodríguez, Mabel Alicia
		
		Roig, María Eugenia 
		
		Rosales, Juan Carlos
		
		Puca, Silvana Mercedes del Milagro 
		
		Saidón, Liliana
		
		Sángari, Antonio Noé 
		
		Micelli, Mónica Lorena
		
		Torres, Germán Ariel
		
		Villagra, Celia Elizabeth
		
		Villarreal Cantizana, Claudia
		
		Viola, Fernanda Beatriz 
		
		Yazlle, Jorge Fernando
	\end{multicols} 
\end{center}

\begin{center}
	\sffamily \Large \bfseries Maquetación
	
	\large \normalfont
	Martinez, Gonzalo Matias
\end{center}

\begin{center}
	\sffamily \Large \bfseries Diseño e identidad
	
	\large \normalfont
	Díaz, Aldana Lucía
\end{center}

\vspace*{2.5mm}

\begin{center}
	\scshape \Large Universidad Nacional de Salta
	
	\normalfont \normalsize
	Av. Bolivia 5150 -- Salta Capital
\end{center}

\vfill

\begin{center}
	\footnotesize
	Queda prohibida la reproducción total o parcial del texto de la presente obra en cualquiera de sus formas, electrónica o mecánica, sin el consentimiento previo y escrito del autor.
\end{center}
	\clearpage
	\begin{center}
	\LARGE \bfseries \sffamily Autoridades
\end{center}

\begin{center}
	\sffamily \Large \bfseries Universidad Nacional de Salta

	\bigskip
	
	\sffamily \large \bfseries Rector
	
	\normalfont \large
	Ing. Daniel HOYOS
	
	\bigskip
	
	\sffamily \large \bfseries Vicerrector
	
	\normalfont \large
	Cr. Nicolás INNAMORATO
	
	\bigskip
	
	\sffamily \large \bfseries Secretario de Extensión 
	
	\normalfont \large
	Lic. Rubén Emilio CORREA
\end{center}

\vspace*{5mm}

\begin{center}
	\sffamily \Large \bfseries Facultad de Ciencias Exactas
	
	\bigskip
	
	\sffamily \large \bfseries Decano
	
	\normalfont \large
	Mg. Gustavo Daniel GIL
	
	\bigskip
	
	\sffamily \large \bfseries Vicedecana
	
	\normalfont \large
	Dra. María Rita MARTEARENA
\end{center}

\vspace*{5mm}

\begin{center}
	\sffamily \Large \bfseries Departamento de Matemática
	
	\bigskip
	
	\sffamily \large \bfseries Directora
	
	\normalfont \large
	Lic. María Cristina AHUMADA
	
	\bigskip
	
	\sffamily \large \bfseries Secretario
	
	\normalfont \large
	Prof. Juan Pablo DIOLI
	
	\bigskip
	
	\sffamily \large \bfseries Prosecretario
	
	\normalfont \large
	Mg. Héctor Nicolás FUNES
\end{center}
	\clearpage
	\begin{center}
	\LARGE \bfseries \sffamily Equipo Editorial JEM
\end{center}

\begin{center}
	\sffamily \Large \bfseries Coordinación General
	
	\bigskip
	
	\normalfont \large
	Prof. Silvia Mabel Baspiñeiro
	
	Prof. Blanca Azucena Formeliano
	
	Prof. lvone Anahí Patagua
	
	Prof. Antonio Noé Sángari
\end{center}

\vspace*{2.5mm}

\begin{center}
	\sffamily \Large \bfseries Plataforma
	
	\bigskip
	
	\sffamily \large \bfseries Open Journal System (OJS)
	
	\normalfont \large
	Lic. Ramiro Guzmán González
	
	\bigskip
	
	\sffamily \large \bfseries Diseño de Identidad
	
	\normalfont \large
	Lic. Aldana Lucía Díaz
	
	\bigskip
	
	\sffamily \large \bfseries Maquetación 
	
	\normalfont \large
	Gonzalo Matias Martinez
\end{center}

\vspace*{5mm}

\begin{center}
	\sffamily \Large \bfseries Declaraciones de Interés Educativo y Avales Institucionales
	
	\bigskip
	
	\normalfont \normalsize \scshape
	Rectorado
	
	Universidad Nacional de Salta
	
	RES. R. N° 937/2022
	
	\bigskip
	
	Facultad de Ciencias Exactas 
	
	Universidad Nacional de Salta
	
	RESD-EXA N° 379/2022
	
	\bigskip
	
	Ministerio de Educación, Cultura, Ciencia 
	
	y Tecnología de la Provincia de Salta
	
	RES. N° 128/2022
	
\end{center}
	\clearpage
	\section*{\centering \sffamily Prólogo}

Las VII Jornadas de Enseñanza de la Matemática híbridas, realizadas en su parte presencial en las instalaciones de la Universidad Nacional de Salta del 24 de julio al 2 de agosto de 2023, representaron un hito en el continuo esfuerzo por mejorar la calidad de la educación matemática en nuestra región. Durante esos días, docentes, investigadores y estudiantes se congregaron con el objetivo de reflexionar sobre sus prácticas docentes, compartiendo experiencias y conocimientos que, sin duda, contribuirán al enriquecimiento de la comunidad educativa.

Las Jornadas ofrecieron un variado programa de actividades que incluyó conferencias magistrales, talleres prácticos y comunicaciones breves. Entre las actividades destacadas, se encuentran las conferencias de Valeria Borsani, que abordó \emph{El pasaje de la aritmética al álgebra en los primeros años de la escuela media}, la de Andrés Rieznik, que abordó \emph{Discalculia, un capítulo olvidado de la neuropsicología}, y el de Daniela Reyes, \emph{Empoderamiento docente, ¿por qué pensar en ello?} y cinco talleres sobre enseñanza de la matemática. La participación activa de los asistentes reflejó el interés y el compromiso con la mejora continua de la enseñanza de la matemática.

Queremos expresar nuestro más sincero agradecimiento a todos los participantes, ponentes y talleristas que, con su dedicación y entusiasmo, hicieron posible el éxito de este evento. Asimismo, extendemos nuestro reconocimiento a la Universidad Nacional de Salta y a todas las instituciones que nos brindaron su apoyo.

Las presentes Memorias recogen las ideas, debates y conclusiones surgidas durante las Jornadas, sirviendo como un valioso recurso para todos aquellos interesados en la enseñanza de la matemática. Esperamos que este compendio inspire a los lectores a continuar explorando y aplicando nuevas estrategias pedagógicas en sus aulas.

Con la mirada puesta en el futuro, confiamos en que las próximas ediciones de las Jornadas de Enseñanza de la Matemática seguirán siendo un espacio de encuentro, reflexión e innovación para nuestra comunidad.

\vspace{5cm}

\begin{center}
	\large Prof. Antonio Noé Sángari
	
	\normalfont Coordinador de las VII Jornadas de Enseñanza de la Matemática
\end{center}
	\clearpage
	\thispagestyle{portadapage}
\tableofcontents
	\clearpage
	\thispagestyle{portadapage}
\renewcommand{\articulotipo}{}

\vspace*{\fill}
\begin{center}
	\sffamily \bfseries \LARGE I
	
	\Huge Talleres
\end{center}
\vspace*{\fill}
\addcontentsline{toc}{part}{I \hspace*{1em} Talleres}
	\clearpage
	%-------------------------------------------------------------------------
% INFORMACIÓN DEL ARTÍCULO
\thispagestyle{portadapage}
\setcounter{subsection}{0}
\setcounter{subsubsection}{0}
\setcounter{actividad}{0}
\setcounter{actividad_previa}{0}
\setcounter{actividad_entre}{0}
\renewcommand{\articulotipo}{Taller}
\renewcommand{\articulotitulo}{Preparación de estudiantes para las Olimpiadas Matemáticas: Técnicas de preparación}
\renewcommand{\articulotitulocorto}{Preparación de estudiantes para las Olimpiadas Matemáticas: Técnicas de preparación}
\section{\articulotitulo}
\desctotoc{Sángari, A. N.; Flores Rocha, V.; Ghiglia, N. N.; Coria, S. E.}

\noindent\rule{\linewidth}{2pt}

\vspace{0.25cm}

\begin{flushright}
	{\Large \scshape
		Antonio Noé Sángari
	}\\
	{\large \itshape
		Universidad Nacional de Salta
	}\\
	{\ttfamily \small
		diamantecinthia@gmail.com
	}\\
	\vspace{1em}
	{\Large \scshape
		Veronica Flores Rocha
	}\\
	{\large \itshape
		Universidad Nacional de Salta
	}\\
	\vspace{1em}
	{\Large \scshape
		Nadia Noel Ghiglia
	}\\
	{\large \itshape
		Universidad Nacional de Salta
	}\\
	\vspace{1em}
	{\Large \scshape
		Silvia Ester Coria
	}\\
	{\large \itshape
		Universidad Nacional de Salta
	}
\end{flushright}

\vspace{0.5cm}

\begin{center}
	\begin{minipage}{0.75\linewidth} \small
		\textsc{Resumen}. ~
		Este curso-taller está diseñado para capacitar a docentes de matemáticas de secundaria en técnicas de demostración matemática, con el objetivo de preparar a sus estudiantes para participar en las Olimpiadas Matemáticas. Se enfocará en el desarrollo de habilidades de razonamiento lógico y deductivo, la enseñanza de estrategias didácticas específicas y la aplicación de estos conocimientos en la práctica. Los participantes tendrán acceso a materiales teóricos, actividades prácticas y herramientas de evaluación para asegurar un aprendizaje integral.
	\end{minipage}
	
	\vspace{0.5em}
	
	\begin{minipage}{0.75\linewidth} \small
		\textsc{Palabras clave} --- Curso-taller, Matemáticas, Olimpiadas Matemáticas, Razonamiento lógico, Razonamiento deductivo.
	\end{minipage}
\end{center}
%-------------------------------------------------------------------------

\subsection{Introducción}

\subsubsection{Sobre el taller}

El curso-taller «Preparación de Estudiantes para las Olimpiadas Matemáticas: Técnicas de Demostración» está diseñado específicamente para docentes de matemáticas de nivel secundario. Este programa tiene como objetivo principal dotar a los profesores de las herramientas y estrategias necesarias para enseñar a sus estudiantes a construir demostraciones matemáticas rigurosas. La habilidad de realizar demostraciones precisas es un aspecto crucial en las Olimpiadas Matemáticas, con un énfasis particular en la geometría básica.

En el contexto de las Olimpiadas Matemáticas, tanto a nivel nacional como internacional, la capacidad de los estudiantes para presentar demostraciones claras y lógicas es fundamental para su éxito. Por esta razón, el curso-taller pone un énfasis especial en proporcionar a los docentes un sólido soporte bibliográfico a través de la Red Olímpica. Esta red ofrece una guía valiosa para las demostraciones esperadas en las competiciones, ayudando a los estudiantes a prepararse adecuadamente y a competir con confianza en los diferentes niveles de las Olimpiadas.

Con este curso-taller, se espera que los docentes adquieran un mayor dominio en la enseñanza de técnicas de demostración y, a su vez, transmitan estos conocimientos a sus alumnos, potenciando así sus habilidades matemáticas y su desempeño en las competiciones.

\subsubsection{Importancia de la enseñanza de la demostración}

Destacamos que en el contexto de la Olimpiada Matemática Argentina, la demostración ocupa un lugar central y fundamental en la resolución de problemas matemáticos. A diferencia de las respuestas simples o los cálculos rápidos, una demostración matemática exige un razonamiento riguroso y lógico.

La demostración es esencial en la Olimpiada Matemática porque asegura la validez de las soluciones, desarrolla habilidades críticas y analíticas, promueve una comunicación efectiva, profundiza la comprensión matemática y permite una evaluación justa y completa de las capacidades
de los estudiantes.

\subsection{Contenidos}

\subsubsection{Módulo 1: Definición y tipos de demostraciones.}
\begin{itemize}
	\item Elementos de una demostración rigurosa.
	\item Ejemplos históricos y su importancia.
	\item Ejemplos de certámenes nacionales e internacionales de la Olimpiada Matemática Argentina, tanto en soporte bibliográfico como en publicación en linea.
	\item Estrategias didácticas para enseñar estos conceptos.
\end{itemize}

\subsubsection{Módulo 2: Técnicas Básicas de Demostración}
\begin{itemize}
	\item Demostraciones directas.
	\item Demostraciones por contradicción.
	\item Demostraciones por contraposición.
	\item Ejercicios prácticos y resolución de problemas.
	\item Métodos para enseñar estas técnicas a los estudiantes.
\end{itemize}

\subsubsection{Módulo 3: Geometría}
\begin{itemize}
	\item Postulados y teoremas fundamentales.
	\item Construcción de demostraciones geométricas.
	\item Aplicación de teoremas en problemas de Olimpiadas en certámenes regionales, nacionales e internacionales.
	\item Actividades prácticas.
\end{itemize}

\subsection{Requisitos previos}

Docentes de matemáticas de nivel secundario y estudiantes avanzados de profesorados de matemática, interesados en preparar a sus estudiantes para participar en Olimpiadas Matemáticas.

\subsection{Objetivos}
\begin{enumerate}
	\item Capacitar a los docentes en técnicas de demostración matemática.
	\item Fomentar habilidades de razonamiento lógico y deductivo en sus estudiantes.
	\item Proveer estrategias didácticas para enseñar demostraciones.
	\item Preparar a los docentes para guiar a sus estudiantes en competiciones matemáticas regionales, nacionales e internacionales.
\end{enumerate}

\subsection{Actividades}

\subsubsection{Actividades previas}

Antes del inicio formal del curso-taller, los participantes tendrán acceso a un curso en la plataforma Moodle de la Facultad de Ciencias Exactas. En esta plataforma, se adjuntarán extractos esenciales de elementos de una prueba en matemática, especialmente del texto de \textcite{margaris1968} sobre lógica formal como soporte básico y del clásico de \textcite{polya1945} sobre resolución de problemas y organización de demostraciones. Además, se proporcionarán videos explicativos que cubren demostraciones básicas para familiarizar a los docentes con los conceptos fundamentales. El programa y el cronograma detallado de actividades del curso estarán disponibles en la plataforma, permitiendo a los participantes planificar su tiempo de estudio. Para asegurar que los docentes asimilen los conceptos preliminares, se implementará un cuestionario de autoevaluación, el cual los participantes deberán completar antes de la primera sesión sincrónica.

\subsubsection{Primeras hora y media sincrónicas}

La primera sesión sincrónica comenzará con una breve presentación del taller, estableciendo los objetivos y la estructura del curso. Se introducirán las ideas básicas de los elementos de una demostración, proporcionando una visión general de los conceptos fundamentales que serán abordados en profundidad más adelante. Se solicitará la participación activa de los cursantes, quienes serán invitados a compartir sus perspectivas y experiencias sobre la enseñanza de demostraciones matemáticas. Esta discusión inicial ayudará a identificar los conocimientos previos de los docentes y a ajustar el enfoque del taller según sus necesidades.

\subsubsection{Primeras tres horas entre clases}

Estas horas estarán
dedicadas a la preparación de las próximas clases sincrónicas centradas en el módulo 2, que aborda diversos métodos de demostración y técnicas para enseñarlos. Se pondrán a disposición textos extraídos del libro de \textcite{larson1983} sobre resolución de problemas para estructurar los métodos de demostración. Adicionalmente, se seleccionarán problemas relevantes de los libros de \textcite{araujo2020a, araujo2020b, araujo2020c, fauring2000, fauring2023}, los cuales serán utilizados como ejemplos prácticos durante las sesiones sincrónicas. Los participantes deberán revisar estos materiales y reflexionar sobre su aplicación en el aula, preparando preguntas y comentarios para la próxima sesión.

\subsubsection{Segundas hora y media sincrónicas}

En esta sesión, se abordará el módulo 2 siguiendo una metodología similar a la de las primeras horas sincrónicas. El enfoque estará en las técnicas básicas de demostración y en métodos pedagógicos para enseñar estas técnicas a los estudiantes. Los cursantes serán invitados a explicar el material que se les asignó para lectura durante las primeras tres horas entre clases, fomentando así una discusión profunda y colaborativa. Se debatirán posibles modificaciones a los ejercicios presentados para obtener respuestas más generales y abarcativas, permitiendo a los docentes adaptar los problemas a diferentes niveles de dificultad y contextos de enseñanza.

\subsubsection{Segundas tres horas entre clases}

Durante estas horas, los participantes prepararán las próximas clases sincrónicas que continuarán con el módulo 3. Este módulo se centrará en la justificación de las hipótesis y en el camino lógico necesario para obtener los teoremas más comunes en la geometría elemental. Se proporcionarán materiales adicionales que exploran estos conceptos en profundidad, incluyendo ejemplos prácticos y ejercicios para resolver. Los docentes deberán revisar estos materiales y preparar sus propias explicaciones y preguntas para la próxima sesión sincrónica, asegurando así una comprensión sólida de los temas tratados.

\subsubsection{Terceras hora y media sincrónicas}

La tercera sesión sincrónica continuará con el módulo 3, siguiendo el mismo enfoque interactivo y participativo de las sesiones anteriores. Se profundizará en las técnicas de demostración y resolución de problemas, con un énfasis especial en la búsqueda de primeros principios para fundamentar los resultados obtenidos. Se dedicará la parte final de esta sesión a explicar la evaluación final del curso-taller, aclarando cualquier duda que los participantes puedan tener. Esto asegurará que todos los cursantes estén bien preparados para demostrar su comprensión y aplicación de las técnicas de demostración matemática enseñadas a lo largo del taller.

\subsubsection{Evaluación final}

\begin{itemize}
	\item Desarrollo de recursos educativos para el entrenamiento de estudiantes, centrados en el razonamiento lógico de deducciones geométricas.
	\item Cuestionario final que aborde los principales temas del taller.
	\item Reflexión escrita sobre la experiencia del taller y el aprendizaje obtenido.
\end{itemize}

\subsection{Bibliografía}

\nocite{*}
\printbibliography[keyword={01}]
	\clearpage
	%-------------------------------------------------------------------------
% INFORMACIÓN DEL ARTÍCULO
\thispagestyle{portadapage}
\setcounter{subsection}{0}
\setcounter{subsubsection}{0}
\setcounter{actividad}{0}
\setcounter{actividad_previa}{0}
\setcounter{actividad_entre}{0}
\renewcommand{\articulotipo}{Taller}
\renewcommand{\articulotitulo}{Datos educativos: producción y uso de herramientas de procesamiento de información con asistencia de la IA para la toma de decisiones}
\phantomsection
\stepcounter{section}
\addcontentsline{toc}{section}{\protect\numberline{\thesection} \articulotitulo}
\desctotoc{Villagra, C. E.; Miguez, I. H; Acosta, R. N.; Chorolque, E. M.}

\begin{center}
	\setstretch{1.5}
	{\Huge \scshape 
		\articulotitulo
	}
\end{center}

\noindent\rule{\linewidth}{2pt}

\vspace{0.25cm}

\begin{flushright}
	{\Large \scshape
		Fernando Jorge Bifano
	}\\
	{\large \itshape
		Facultad de Ciencias Exactas y Naturales — Universidad de Buenos Aires
	}\\
	{\ttfamily \small
		fjbifano@ccpems.exactas.uba.ar
	}\\
	\vspace{1em}
	{\Large \scshape
		Pablo Fabián Carranza
	}\\
	{\large \itshape
		Universidad Nacional de Río Negro
	}\\
	\vspace{1em}
\end{flushright}

\vspace{0.5cm}

\begin{center}
	\begin{minipage}{0.75\linewidth} \small
		\textsc{Resumen}. ~
		La generación y uso de datos es actualmente uno de los nuevos emergentes que impactan en las diferentes esferas de la sociedad y por tanto, la educación se ve tensionada por las demandas que ello conlleva. En ese sentido, surgen nuevas vacancias para la formación docente y se vuelve necesario ofrecer instancias que permitan a quienes enseñan contar con herramientas para la comprensión, uso y toma de decisiones fundadas y sostenidas en la interpretación de la información. Este taller ofrece la posibilidad de familiarizarse con herramientas básicas que brindan los softwares para el procesamiento de la información y que con la asistencia de las IA se ven potenciados en sus aplicaciones. A partir del trabajo con distintos tipos de bases de datos educativos, propondremos un recorrido que aborda algunas herramientas de representación gráfica de la información en lenguaje de programación Python asistido por medio de inteligencias artificiales tales como ChatGPT y Gemini. Se analizarán las potencialidades de estrategias tales como clustering y árboles de decisión para la caracterización de grupos de estudiantes.
	\end{minipage}
	
	\vspace{0.5em}
	
	\begin{minipage}{0.75\linewidth} \small
		\textsc{Palabras clave} --- Datos educativos, IA, Python, Procesamiento de la información.
	\end{minipage}
\end{center}
%-------------------------------------------------------------------------

\subsection{Introducción}

El análisis de datos es una disciplina que ha visto incrementado su desarrollo recientemente por varios factores vinculados entre sí. Entre ellos destacamos:

\begin{itemize}
	\item Disponibilidad de herramientas de procesamiento de datos en computadoras de escritorio.
	\item Crecimiento de librerías de libre acceso en lenguaje Python.
	\item Desarrollo de nuevos métodos de análisis de datos, llegando incluso a herramientas vinculadas a la inteligencia artificial.
	\item Existencia de volúmenes de datos.
\end{itemize}

Estos elementos combinados permiten la aplicación de herramientas que facilitan la comprensión de fenómenos complejos y la toma de decisiones basadas en información.

Las instituciones educativas no escapan a estos avances. En efecto, los datos disponibles o factibles de ser obtenidos permiten una mejor interpretación de fenómenos del ecosistema educativo y el desarrollo de acciones basadas en información. A modo de ejemplo citamos:

\begin{itemize}
	\item análisis de desempeño integral de estudiantes
	\item relaciones entre estrategias de estudio y calificaciones
	\item relaciones entre proyectos de vida, resiliencia y retención universitaria
	\item sistemas de alerta temprana de abandono
	\item caracterización de grupos de estudiantes por técnicas de clustering
\end{itemize}

El análisis de datos, entonces, resulta una herramienta que permite extraer información para la comprensión de fenómenos que acontecen en la institución educativa y así tomar decisiones basadas en información precisa y contextualizada.

Estas herramientas pueden ser utilizadas tanto por los equipos de conducción como por los docentes en la aulas; en contextos presenciales o virtuales. En esta propuesta nos centraremos en algunas herramientas de análisis de datos que resultan de interés a docentes tanto sea al interior de una cátedra como para el análisis interdisciplinario del desempeño de estudiantes. Las mismas son también de gran interés para la gestión de la institución, facilitando la comprensión de la dinámica de la misma y brindando fundamentos para la toma de decisiones.

\subsection{Contenidos}

Se abordarán algunas herramientas de representación gráfica de la información en lenguaje de programación Python asistido por medio de inteligencias artificiales tales como ChatGPT y Gemini. Se analizarán las potencialidades de herramientas tales como clustering y árboles de decisión para la caracterización de grupos de estudiantes. Más precisamente se proponen los siguientes módulos:

\subsubsection{Actividades previas}

\begin{enumerate}[a)]
	\item Introducción a lenguaje de programación Python en Jupyter	Notebook
	
	El análisis de datos se realizará en el entorno llamado Jupyter Notebook, herramienta accesible desde el software Anaconda (gratuito).
	
	\item Abrir una base de datos en formatos xlsx o csv en Jupyter	Notebook con Python
	
	Los datos serán analizados en python, por lo que el primer paso será subir la base de datos al entorno Jupyter Notebook para poder tratar la base de datos	en el lenguaje python.
\end{enumerate}

\subsubsection{Módulo 1}

En este módulo se abordarán los conceptos básicos de programación en Python en Jupyter Notebook para un primer tratamiento de datos educativos. En esta primera etapa, se trabajará con datos numéricos reales anonimizados de alumnos de escuelas secundarias de Argentina. Entre otros, serán tratados temas tales como:

\begin{enumerate}[a)]
	\item primera exploración de la base de datos
	\item representación gráfica de variables
	\item correlaciones
	\item normalización o estandarización
	\item clustering por método k-nn
	\item árboles de decisión
\end{enumerate}

Todos los métodos serán considerados como herramientas de síntesis para la extracción de información de los datos. En todas las instancias se accederá al uso de IA generativas como asistentes a la programación en Python.

En este módulo se debatirá sobre las condiciones del trabajo final del taller, consistente en una presentación por dupla de un análisis de datos propios. Más precisamente, se consagrarán los últimos minutos a establecer las características que deben reunir los datos para un análisis acorde a las posibilidades del taller.

\subsubsection{Módulo 2}

En este módulo se abordará el tratamiento de datos categóricos en escala de likert, muy frecuentes en contextos educativos. Se tratarán en principio con los mismos métodos abordados en el módulo 1. Se integrarán también con datos de origen numérico para posibles extracciones de información de tales cruzamientos.

Se dedicarán los últimos minutos a terminar de definir las características de la presentación que cada dupla hará sobre el análisis de datos propios como trabajo final.

\subsubsection{Módulo 3}

Este módulo se consagrará enteramente a las presentaciones que cada dupla haya podido realizar sobre el análisis de datos educativos propios. Se realizarán también devoluciones y comentarios de parte de todos los participantes a las diferentes presentaciones.

\subsection{Requisitos previos}

Este taller está destinado a docentes que enseñan matemáticas en los niveles medio y superior. Se requiere conocimientos mínimos en el área de la estadística descriptiva. Así mismo resulta deseable que los asistentes cuenten con cierta familiaridad en relación con el uso de herramientas tecnológicas tales como Excel. El manejo de lenguaje de programación previo no es excluyente y se espera que sea una de las capacidades a desarrollar en los asistentes.

\subsection{Objetivos}

En este taller, nos proponemos alcanzar los siguientes objetivos:
\begin{itemize}
	\item Favorecer la reflexión crítica sobre la producción y el uso de datos 	educativos para la toma de decisiones fundamentadas.
	\item Introducir en el uso de los rudimentos básicos de la programación en Python para el análisis de datos.
	\item Utilizar herramientas tecnológicas, especialmente con la asistencia de la IA, para potenciar el análisis de la información educativa para la toma de decisiones.
	\item Desarrollar habilidades básicas en el tratamiento de bases de datos para el análisis estadístico.
\end{itemize}

\subsection{Actividades}

\subsubsection{Actividades previas}

Para poder participar plenamente del taller, les proponemos a los cursantes que como actividades previas se familiaricen con algunos softwares y/o herramientas que utilizaremos durante el desarrollo del mismo. Específicamente les proponemos:

\begin{itemize}
	\item Para una introducción al tema sugerimos seguir las propuestas detalladas en los siguientes vídeos:
	
	Descargar Anaconda: \url{https://www.anaconda.com/download}
	
	Vídeos de primeros pasos de Python en Jupyter Notebook:
	\begin{itemize}
		\item \url{https://www.youtube.com/watch?v=81-AyuxjdBo}
		\item \url{https://www.youtube.com/watch?v=cp5gUF1D0nQ\&list=PLMUoURdFUxkkHw8tVweJi8YAZcZHDjYou}
	\end{itemize}
	
	\item Para una introducción a la etapa de acceso a una base de datos desde Jupyter notebook, los participantes pueden acceder al siguiente vídeo:
	\begin{itemize}
		\item \url{https://www.youtube.com/watch?v=dbEQtzObsQw}
	\end{itemize}
	
	Los enlaces no son excluyentes, los participantes encontrarán muchos otros en la web.
\end{itemize}

\subsubsection{Primera hora y media sincrónica}\label{primera-hora-y-media-sincronica}

A lo largo de la primera etapa sincrónica, le propondremos a los cursantes, el análisis de una base de datos anonimizada basada en resultados de rendimiento académico (módulo 1) —provista por los profesores a cargo del taller— que se centrará en aspectos tales como: representación gráfica de datos, correlaciones, clustering, entre otros.
\begin{itemize}
	\item Exploración de base de datos.
	\item Representación gráfica de distribuciones de variables.
	\item Establecimiento de posibles relaciones entre variables.
	\item Utilización de métodos de machine learning tales como clustering y árboles de decisión.
\end{itemize}

\subsubsection{Primeras tres horas entre clases}\label{primeras-tres-horas-entre-clases}

Como parte de la primera sección entre clases, los estudiantes en grupos deberán producir un breve informe con los principales hallazgos obtenidos a partir del análisis estadístico efectuado en la sección \ref{primera-hora-y-media-sincronica}. Las características del mismo son:
\begin{itemize}
	\item Una extensión máxima de 5 carillas incluyendo gráfico y/o anexos.
	\item Expresar clara y fundadamente, los principales resultados que surgen del análisis estadístico de la base de datos estudiada.
\end{itemize}

\subsubsection{Segundas hora y media sincrónicas}\label{segundas-hora-y-media-sincronicas}

La primera parte de la sesión estará abierta a compartir algunos de los principales resultados producidos en los informes grupales como fruto de la etapa \ref{primeras-tres-horas-entre-clases}. En una segunda parte, análogamente a lo descrito en la sección \ref{primera-hora-y-media-sincronica}, los profesores del taller ofreceremos a los cursantes una base de datos educativos relacionados con el abandono escolar para su análisis (módulo 2).

\subsubsection{Segundas tres horas entre clases}\label{segundas-tres-horas-entre-clases}

Este espacio es análogo al descrito en la sección \ref{primeras-tres-horas-entre-clases}. La diferencia fundamental radica en que se les propondrá a los cursantes del taller, la búsqueda y/o elaboración de información que pueda ser un insumo para la construcción de una base de datos propia que les permita hacer un estudio de alguna problemática relacionada con las instituciones escolares en la que se desempeñan. Esto será el insumo fundamental para la producción final.

\subsubsection{Terceras hora y media sincrónicas}

Este espacio es análogo al descrito en la sección \ref{segundas-hora-y-media-sincronicas}. A diferencia de la etapa referenciada, dedicaremos una primera parte de la sesión a que los estudiantes compartan los datos recogidos en la etapa \ref{segundas-tres-horas-entre-clases}. a los fines de discutir la viabilidad de su estudio para la producción final (módulo 3).

\subsubsection{Evaluación final}

Como parte de la evaluación final del curso, propondremos a los cursantes las siguientes actividades:
\begin{itemize}
	\item Elaboración de un proyecto final aplicando los conocimientos adquiridos. Más precisiones sobre el mismo se ofrecerán oportunamente a través de la plataforma de la facultad.
	\item Cuestionario final que aborde los principales temas del taller a modo de autoevaluación del taller.
	\item Reflexión escrita sobre la experiencia del taller y el aprendizaje obtenido dando cuenta de los principales obstáculos superados a lo largo de la experiencia.
\end{itemize}

\subsection{Bibliografía}

\nocite{*}
\printbibliography[keyword={02}]
	\clearpage
	%-------------------------------------------------------------------------
% INFORMACIÓN DEL ARTÍCULO
\thispagestyle{portadapage}
\setcounter{subsection}{0}
\setcounter{subsubsection}{0}
\setcounter{actividad}{0}
\setcounter{actividad_previa}{0}
\setcounter{actividad_entre}{0}
\renewcommand{\articulotipo}{Taller}
\renewcommand{\articulotitulo}{Autovalores, autovectores y transformaciones: una mirada geométrica}
\phantomsection
\stepcounter{section}
\addcontentsline{toc}{section}{\protect\numberline{\thesection} \articulotitulo}
\desctotoc{Chañi, M. D.; Romero, S. N.; Gutierrez, M. D.; Velasquez, N. A. M.}

\begin{center}
	\setstretch{1.5}
	{\Huge \scshape 
		\articulotitulo
	}
\end{center}

\noindent\rule{\linewidth}{2pt}

\vspace{0.25cm}

\begin{flushright}
	{\Large \scshape
		Marcos Darío Chañi
	}\\
	{\large \itshape
		Universidad Nacional de Salta
	}\\
	{\ttfamily \small
		mat.marcos2.0@gmail.com
	}\\
	\vspace{1em}
	{\Large \scshape
		Silvia Noemí Romero
	}\\
	{\large \itshape
		Universidad Nacional de Salta
	}\\
	\vspace{1em}
	{\Large \scshape
		Marcela Daniela Gutierrez
	}\\
	{\large \itshape
		Universidad Nacional de Salta
	}\\
	\vspace{1em}
	{\Large \scshape
		Noelia Adriana Melisa Velasquez
	}\\
	{\large \itshape
		Universidad Nacional de Salta
	}\\
\end{flushright}

\vspace{0.5cm}

\begin{center}
	\begin{minipage}{0.75\linewidth} \small
		\textsc{Resumen}. ~
		Dentro de la enseñanza del matemática, más precisamente el álgebra lineal uno de los tópicos fundamentales el concepto de «Transformación Lineal», y a veces por razones de tiempo no se puede profundizar en algunos temas complementarios. Es así que se propone en este taller resignificar los temas de «Transformaciones Lineales, autovalor y autovector» entre otros, para abordar un nuevo tema que son las «Transformaciones ortogonales», así mediante la ayuda del álgebra lineal se espera que los alumnos logren interpretar estos conceptos de manera gráfica y geométrica, para que no sean sólo un conjunto de conocimientos abstractos. Este trabajo ofrece una alternativa de enseñanza para los conceptos de diferentes transformaciones. El taller es un espacio para realizar gráficos dinámicos con GeoGebra, acompañado de situaciones de análisis y de debate, con la intención de realizar actividades matemáticas mediadas por un instrumento tecnológico donde el eje central es el concepto matemático, y de esta manera superar las dinámicas tradicionales. Así se propone una nueva forma de emprender las prácticas educativas contextualizas en los actuales paradigmas educativos.
	\end{minipage}\\
	
	\vspace{0.5em}
	
	\begin{minipage}{0.75\linewidth} \small
		\textsc{Palabras clave} --- GeoGebra, Álgebra, Trasformaciones Ortogonales, Matrices, Autovalor, Autovector.
	\end{minipage}
\end{center}
%-------------------------------------------------------------------------

\subsection{Introducción}

\subsubsection{Importancia del Taller}

Dentro de la enseñanza del matemática, más precisamente el álgebra lineal uno de los tópicos fundamentales el concepto de <<Transformación Lineal>>, y a veces por razones de tiempo no se puede profundizar en algunos temas complementarios. Es así que se propone en este taller resignificar los temas de «Transformaciones Lineales, autovalor y autovector» entre otros, para abordar un nuevo tema que son las <<Transformaciones ortogonales>>, así mediante la ayuda del álgebra lineal se espera que los alumnos logren interpretar estos conceptos de manera gráfica y geométrica, para que no sean sólo un conjunto de conocimientos abstractos. Este trabajo ofrece una alternativa de enseñanza para los conceptos de diferentes transformaciones. El taller es un espacio para realizar gráficos dinámicos con GeoGebra, acompañado de situaciones de análisis y de debate, con la intención de realizar actividades matemáticas mediadas por un instrumento tecnológico donde el eje central es el concepto matemático, y de esta manera superar las dinámicas tradicionales. Así se propone una nueva forma de emprender las prácticas educativas contextualizas en los actuales paradigmas educativos.

\subsubsection{Fundamentos}

En relación al marco pedagógico se va a trabajar a través del aprendizaje secuencial y colaborativo, mediante un “conocimiento espiralado”, es decir el alumno construye y/o afirma conceptos por medio de conceptos ya aprendidos, de manera que realiza un reordenamiento de sus ideas y los confirma. Es por ello que se propone un taller, ya que las actividades pueden realizarse en forma individual o colectiva, la evaluación de las actividades es realizada por los integrantes, de esta manera se ubica al alumno para que sea crítico tanto con su trabajo como el de sus pares. 

Se elaborará una guía para los 3 encuentros, en las cuales se proponen actividades que se realizarán en forma grupal o individual en forma alternada. Las actividades se llevarán a cabo en tres instancias las cuales habrá:
\begin{itemize}
	\item Un desarrollo matemático en función a lo que se pide en las guías.
	\item Un trabajo en la máquina mediante el uso del Software GeoGebra.
	\item Una puesta en común para debatir sobre lo realizado en las instancias anteriores.
\end{itemize}

Los nuevos escenarios educativos brindan la posibilidad de incorporar recursos vinculados con las TIC, así también cada día aumenta la
demanda de la formación continua en relación a las nuevas corrientes tecnológicas-matemáticas por parte de los docentes, esto conlleva
a generar espacios de aprendizaje en donde puedan elaborar y llevar a cabo propuestas educativas que involucren las TIC en la enseñanza de la matemática. 

El presente escrito describe un taller de actividades mediante el uso del software GeoGebra, con la intención de contribuir a la formación profesional continua en el campo de la matemática y su enseñanza. La elección del software GeoGebra, se debe no solo a sus potencialidades matemáticas y didácticas, sino también a que es de uso libre.

Hoy en día el desafío de enseñar matemática implica que los docentes no solo superen algunas tradiciones que provienen de nuestras biografías escolares y formativas, sino también el replantear el uso de las tecnologías en el proceso de enseñanza y aprendizaje. A partir de esto se realiza una propuesta para estudiantes avanzados y docentes de matemática, con la intención de realizar actividades matemáticas mediadas por un instrumento tecnológico con eje central el concepto matemático, y de esta manera superar las dinámicas tradicionales.

Uno de los aspectos del taller es la visualización, en relación a ella distintos investigadores \textcite{castro1997,cantoral2002} acuerdan que al realizar una actividad de visualización se requiere la utilización de nociones y conceptos matemáticos para poder ser interpretada de forma adecuada. Es decir, la capacidad de visualizar un concepto matemático o un problema requiere de la habilidad para interpretar y entender la información de tal, manipularlo mentalmente y así expresarlo en un soporte material. Así también cuando se usa la representación gráfica de un concepto matemático como herramienta para interpretar otros o resolver problemas, la visualización no es el fin en si mismo sino un medio para llegar a la compresión de propiedades y de relaciones entre distintos conceptos. Según \textcite[103]{castro1997}, dominar un concepto matemático consiste en conocer sus principales representaciones y el significado de cada una de ellas, así como operar con las reglas internas de cada sistema y en convertir o traducir unas representaciones en otras, detectando qué sistema es más ventajoso para trabajar cono determinadas propiedades. 

Debe remarcarse que GeoGebra permite otros modos de hacer en la clase ya que los participantes pueden enfocar su atención en procesos de análisis, que permiten la modelización, la toma de decisiones, el razonamiento y la resolución de situaciones, tal como lo plantean \textcite{arcavi2000} (citado en Santos Trigo, 2007, p. 39), los ambientes dinámicos no sólo permiten a los estudiantes construir figuras con ciertas propiedades y visualizarlas, sino que también les permite transformar esas construcciones en tiempo real. Este dinamismo puede contribuir en la formación de hábitos para transformar (mentalmente o por medio de una herramienta) una instancia particular, para estudiar variaciones, invariantes visuales y posiblemente proveer bases intuitivas para justificaciones formales de conjeturas y proposiciones.

\subsection{Contenidos}

Interpretación geométrica de transformaciones lineales, autovalor y autovector de una transformación lineal. Transformaciones ortogonales. 

\subsection{Requisitos Previos}

Conceptos básicos de un curso de álgebra lineal. Especialmente Transformaciones lineales. Autovalor y autovector.

\subsection{Objetivos}
\begin{itemize}
	\item Valorizar el concepto de transformación lineal, autovalor y autovector a través de su interpretación geométrica.
	\item Facilitar el aprendizaje de las diferentes transformaciones ortogonales, a través de la visualización gráfica. mediante el Software GeoGebra.
	\item Realizar representaciones gráficas dinámicas de los conceptos descritos anteriormente.
	\item Generar en los participantes nuevas ideas de cómo utilizar los recursos vistos, en su práctica docente.
\end{itemize}

\subsection{Actividades}

\subsubsection{Actividades previas}

Lectura del material bibliográfico sobre el tema <<Transformaciones lineales, autovalores y autovectores>> sumado a los ejercicios de identificar si una transformación es lineal o no, y sobre autovalores y autovectores. Se espera que mediante las actividades los participantes refresquen dichos contenidos que fueron impartidos en un curso de álgebra lineal.

\subsubsection{Primera hora y media presenciales}\label{subsec:Primeras-dos-sinc-03}

Los cursantes del taller a partir de la <<Definición de Transformación Ortogonal>>, mediante razonamientos lógicos y deductivos deberán encontrar las condiciones para ser una transformación ortogonal y de esta manera hallar las dos maneras de expresar una transformación ortogonal en forma matricial, a las que llamaremos de primera y segunda especie. Cuando los cursantes lo requieran contarán con la guía de los docentes a cargo para poder evacuar dudas o confirmar si sus razonamientos son correctos y poder validarlos.

\subsubsection{Primeras dos horas entre clases}\label{subsec:Primeras-dos-EC-03}

En este espacio se dará material tanto teórico como práctico, para el uso de GeoGebra. En relación al material teórico se propone un resumen de cómo introducir matrices en el software GeoGebra, para luego pasar a las actividades donde deberán introducir las diferentes matrices que se proponen. Los archivos se guardarán, para en forma posterior ser enviados mediante la plataforma para su control por parte de los responsables.

\subsubsection{Segundas dos horas presenciales}

En las segunda dos horas presenciales, los cursantes del taller deberán deducir la expresión matricial de diferentes transformaciones rígidas en $\mathbb{R}^{3}$ es decir en el plano (traslación, simetría central, simetría axial, etc.) con el sustento de un un marco gráfico, para ello deberán hacer uso del concepto de autovalor y autovector de una transformación lineal. Dichas actividades estarán guiadas mediante preguntas y/o observaciones de manera que los participantes del taller puedan razonar y debatir ya sea de manera individual o grupal. En todo el proceso de deducción y discusión los docentes responsables podrán guiar a los participantes para que logren hallar las expresiones matriciales de cada transformación rígida.

\subsubsection{Segundas dos horas entre clases}

Este espacio estará dedicado para que los participantes puedan validar mediante el Software GeoGebra que las transformaciones halladas son correctas o no. Para ello se les pedirá en GeoGebra en la ventada 2D:
\begin{itemize}
	\item Construir una polígono (por ejemplo un triángulo).
	\item Ingresar la matriz de la transformación correspondiente (se aclara que en las primeras horas entre clases, sección \ref{subsec:Primeras-dos-EC-03}, los participantes ya aprendieron a cargar matrices en GeoGebra).
	\item Encontrar mediante la matriz ingresada, los transformados de los vértices de la figura y construir en nuevo polígono con los puntos resultantes de aplicar dicha transformación.
	\item Comprobar, validar o refutar que dicha es la correcta.
\end{itemize}

\subsubsection{Terceras dos horas presenciales}

En las terceras y últimas dos horas presenciales, los cursantes nuevamente deberán deducir la expresión matricial de las transformaciones estudiadas en las dos horas presenciales anteriores, sección \ref{subsec:Primeras-dos-sinc-03}, pero en este caso para el espacio $\mathbb{R}^{3}$. Incentivando el poder de generalización de los participantes. Dichas actividades nuevamente estarán guiadas mediante preguntas y/o observaciones de manera que los participantes puedan generalizar el método encontrado en la clase anterior. Así también en todo el proceso de deducción y discusión los docentes responsables podrán sugerir a los participantes pensar en cómo generalizar las expresiones matriciales de cada transformación rígida.

\subsubsection{Evaluación final}

La evaluación será mediante la entrega de archivos de GeoGebra, para ello a participantes se les dará una transformación en $\mathbb{R}^{3}$ a realizar, la cual deberán realizar la matriz de transformación y luego ingresarla en el Software. En forma posterior se le pedirá que verifiquen que una figura (polígono o poliedro) se transforma de manera correcta mediante el uso de la matriz ingresada.

\subsection{Bibliografía}

\nocite{*}
\printbibliography[keyword={03}]
	\clearpage
	\thispagestyle{portadapage}
\renewcommand{\articulotipo}{}

\vspace*{\fill}
\begin{center}
	\sffamily \bfseries \LARGE II
	
	\Huge Comunicaciones breves
\end{center}
\vspace*{\fill}
\addcontentsline{toc}{part}{II \hspace*{1em} Comunicaciones breves}
	\clearpage
	\includepdf{Trabajos/Contraportada.pdf}

\end{document}